% Auto-generated from index.mdx — edit index.mdx then re-run mdx2tex.mjs
\chapter{前言}
\label{ch:index}

\section*{Python 装饰器 CookBook}
\begin{tcolorbox}[colback=blue!5,colframe=blue!40,title={📘 本册地图}]
本 CookBook 面向中级 Python 开发者，系统讲解装饰器（Decorator）这一元编程核心工具。

\begin{itemize}
\item 第 1 章 · 装饰器基础：闭包、\textbackslash\{\}textbackslash\textbackslash\{\}\{\textbackslash\{\}\}texttt\textbackslash\{\}\{@\textbackslash\{\}\} 语法糖、无参装饰器 - 第 2 章 · 进阶装饰器：带参装饰器、类装饰器、装饰器堆叠 - 附录 A · functools 速查
\end{itemize}
\end{tcolorbox}
\section{为什么要学装饰器}
装饰器是 Python 中最具"Pythonic"风格的语言特性之一。从 Web 框架（Flask 的 \textbackslash\{\}texttt\{@app.route\}、FastAPI 的 \textbackslash\{\}texttt\{@app.get\}）到权限校验、日志记录、缓存、重试，几乎所有横切关注点（cross-cutting concern）都可以用装饰器优雅地表达。

理解装饰器不仅是会写 \textbackslash\{\}texttt\{@xxx\}，更要理解：

\begin{itemize}
\item \textbackslash\{\}textbf\{函数是一等公民\}（first-class citizen）：函数可以作为参数、返回值、赋值给变量。
\item \textbackslash\{\}textbf\{闭包\}（closure）：内部函数记住外部函数的词法环境。
\item \textbackslash\{\}textbf\{语法糖只是函数调用\}：\textbackslash\{\}texttt\{@deco\} 本质上是 \textbackslash\{\}texttt\{f = deco(f)\}。
\end{itemize}
\section{全景图}

\begin{tcolorbox}[colback=gray!5,colframe=gray!40,title={Diagram: mermaid}]
Diagram source kept in \texttt{diagrams/}. Rendered in Nextra/PDF.\end{tcolorbox}

> 上面这张心智图（mindmap）覆盖了本册的全部知识点。建议先通览，再进入第 1 章。

\section{目标读者}
\begin{itemize}
\item 已经会写 Python 函数与类，读过至少一份 Python 项目源码
\item 想系统理解 \textbackslash\{\}texttt\{@decorator\} 的内部机理，而不是只会复制粘贴
\item 目标：能写出可复用、可维护、可调试的装饰器
\end{itemize}
\section{约定}
\begin{itemize}
\item 所有示例基于 \textbackslash\{\}textbf\{Python 3.10+\}（使用 PEP 604 的 \textbackslash\{\}texttt\{X | Y\} 类型注解语法）。
\item 代码块标注文件路径，可直接复制运行。
\item 每章末尾的 \textbackslash\{\}textbf\{Recipes\} 提供任务导向示例；\textbackslash\{\}textbf\{Pitfalls\} 列出常见坑。
\end{itemize}
